% !TEX spellcheck = en_US
% !TEX spellcheck = LaTeX
\documentclass[letterpaper,10pt,english]{article}
\usepackage{%
amsfonts,%
amsmath,%
amssymb,%
amsthm,%
babel,%
bbm,%
%biblatex,%
caption,%
centernot,%
color,%
enumerate,%
%enumitem,%
epsfig,%
epstopdf,%
etex,%
fancybox,%
framed,%
fullpage,%
%geometry,%
graphicx,%
hyperref,%
latexsym,%
mathptmx,%
mathtools,%
multicol,%
paralist,%
pgf,%
pgfplots,%
pgfplotstable,%
pgfpages,%
proof,%
psfrag,%
%subfigure,%
tcolorbox,%
tfrupee,%
tikz,%
times,%
%ulem,%
url,%
xcolor,%
mathpazo,%
}

\definecolor{shadecolor}{gray}{.95}%{rgb}{1,0,0}
\usepackage[margin=1in,top=0.75in]{geometry}
\usepackage[mathscr]{eucal}
\usepgflibrary{shapes}
\usepgfplotslibrary{fillbetween}
\usetikzlibrary{%
arrows,%
backgrounds,%
chains,%
decorations.pathmorphing,% /pgf/decoration/random steps | erste Graphik
decorations.text,% 
matrix,%
positioning,% wg. " of "
fit,%
patterns,%
petri,%
plotmarks,%
scopes,%
shadows,%
shapes.misc,% wg. rounded rectangle
shapes.arrows,%
shapes.callouts,%
shapes%
}

%\pgfplotsset{compat=newest} %<------ Here
\pgfplotsset{compat=1.11} %<------ Or use this one

\theoremstyle{plain}
\newtheorem{thm}{Theorem}[section]
\newtheorem{lem}[thm]{Lemma}
\newtheorem{prop}[thm]{Proposition}
\newtheorem{cor}[thm]{Corollary}
\newtheorem{clm}[thm]{Claim}

\theoremstyle{definition}
\newtheorem{defn}[thm]{Definition}
\newtheorem{conj}[thm]{Conjecture}
\newtheorem{exmp}[thm]{Example}
\newtheorem{axiom}[thm]{Axiom}
\newtheorem{assum}[thm]{Assumption}
\newtheorem{exerc}[thm]{Exercise}
\newtheorem*{defin}{Definition}

\theoremstyle{remark}
\newtheorem{rem}{Remark}
\newtheorem{note}{Note}

\newcommand{\iid}{\emph{i.i.d.}\ }
\newcommand{\Cov}{\operatorname{Cov}}
%\newcommand{\det}{\operatorname{det}}
%\newcommand{\Real}{\mathbb{R}}
\newcommand{\tr}{\operatorname{tr}}
\newcommand{\Var}{\operatorname{Var}}
\newcommand{\cov}{\operatorname{cov}}

%\renewcommand{\proof}[1]{\begin{proof}#1\end{proof}}
\newcommand{\EQ}[1]{\begin{equation*}#1\end{equation*}}
\newcommand{\EQN}[1]{\begin{equation}#1\end{equation}}
\newcommand{\eq}[1]{\begin{align*}#1\end{align*}}
\newcommand{\meq}[2]{\begin{xalignat*}{#1}#2\end{xalignat*}}
\newcommand{\norm}[1]{\left\lVert#1\right\rVert}
\newcommand{\inner}[1]{\left\langle #1\right\rangle}
\newcommand{\abs}[1]{\left\lvert#1\right\rvert}
\newcommand{\expect}[1]{\mathbb{E}\left[{#1}\right]}
\newcommand{\prob}[1]{\mathbb{P}\left[{#1}\right]}
\newcommand{\given}{\; \big\vert \;} 
\newcommand{\set}[1]{\left\{#1\right\}} 
\newcommand{\indicator}[1]{1_{\set{#1}}} 
\newcommand{\Ind}[1]{\mathbbm{1}_{#1}} 
\newcommand{\SetIn}[1]{\mathbbm{1}_{\set{#1}}} 
\newcommand{\red}[1]{\textcolor{red}{#1}} 
\newcommand{\floor}[1]{\left\lfloor#1\right\rfloor}
\newcommand{\ceil}[1]{\left\lceil#1\right\rceil}


\newcommand{\C}{\mathbb{C}}
\newcommand{\D}{\mathbb{D}}
\newcommand{\E}{\mathbb{E}}
\newcommand{\F}{\mathbb{F}}
\newcommand{\N}{\mathbb{N}}
\renewcommand{\P}{\mathbb{P}}
\newcommand{\Q}{\mathbb{Q}}
\newcommand{\R}{\mathbb{R}}
\newcommand{\Z}{\mathbb{Z}}

\newcommand{\bA}{\mathbf{A}}
\newcommand{\bI}{\mathbf{I}}
\newcommand{\bU}{\mathbf{1}}
\newcommand{\bx}{\mathbf{x}}
\newcommand{\bX}{\mathbf{X}}
\newcommand{\bZ}{\mathbf{Z}}


\newcommand{\cA}{\mathcal{A}}
\newcommand{\cB}{\mathcal{B}}
\newcommand{\cC}{\mathcal{C}}
\newcommand{\cF}{\mathcal{F}}
\newcommand{\cG}{\mathcal{G}}
\newcommand{\cM}{\mathcal{M}}
\newcommand{\cN}{\mathcal{N}}
\newcommand{\cP}{\mathcal{P}}
\newcommand{\cT}{\mathcal{T}}
\newcommand{\cX}{\mathcal{X}}
\newcommand{\cY}{\mathcal{Y}}

\newcommand{\sA}{\mathscr{A}}
\newcommand{\sB}{\mathscr{B}}
\newcommand{\sC}{\mathscr{C}}
\newcommand{\sE}{\mathscr{E}}
\newcommand{\sF}{\mathscr{F}}
\newcommand{\sG}{\mathscr{G}}
\newcommand{\sH}{\mathscr{H}}
\newcommand{\sL}{\mathscr{L}}
\newcommand{\sS}{\mathscr{S}}
\newcommand{\sT}{\mathscr{T}}
\newcommand{\sX}{\mathscr{X}}
\newcommand{\sY}{\mathscr{Y}}
\newcommand{\sZ}{\mathscr{Z}}

\newcommand{\tS}{\tilde{S}}
% Debug
\newcommand{\todo}[1]{\begin{color}{blue}{{\bf~[TODO:~#1]}}\end{color}}

% a few handy macros

\renewcommand{\le}{\leqslant}
\renewcommand{\ge}{\geqslant}
\newcommand\matlab{{\sc matlab}}
\newcommand{\goto}{\rightarrow}
\newcommand{\bigo}{{\mathcal O}}
%\newcommand{\half}{\frac{1}{2}}
%\newcommand\implies{\quad\Longrightarrow\quad}
\newcommand\reals{{{\rm l} \kern -.15em {\rm R} }}
\newcommand\complex{{\raisebox{.043ex}{\rule{0.07em}{1.56ex}} \hskip -.35em {\rm C}}}


% macros for matrices/vectors:

% matrix environment for vectors or matrices where elements are centered
\newenvironment{mat}{\left[\begin{array}{ccccccccccccccc}}{\end{array}\right]}
\newcommand\bcm{\begin{mat}}
\newcommand\ecm{\end{mat}}

% matrix environment for vectors or matrices where elements are right justifvied
\newenvironment{rmat}{\left[\begin{array}{rrrrrrrrrrrrr}}{\end{array}\right]}
\newcommand\brm{\begin{rmat}}
\newcommand\erm{\end{rmat}}

% for left brace and a set of choices
%\newenvironment{choices}{\left\{ \begin{array}{ll}}{\end{array}\right.}
\newcommand\when{&\text{if~}}
\newcommand\otherwise{&\text{otherwise}}
% sample usage:
%\delta_{ij} = \begin{choices} 1 \when i=j, \\ 0 \otherwise \end{choices}


% for labeling and referencing equations:
\newcommand{\eql}{\begin{equation}\label}
\newcommand{\eqn}[1]{(\ref{#1})}
% can then do
%\eql{eqnlabel}
%...
%\end{equation}
% and refer to it as equation \eqn{eqnlabel}.
% some useful macros for finite difference methods:
\newcommand\unp{U^{n+1}}
\newcommand\unm{U^{n-1}}
% for chemical reactions:
\newcommand{\react}[1]{\stackrel{K_{#1}}{\rightarrow}}
\newcommand{\reactb}[2]{\stackrel{K_{#1}}{~\stackrel{\rightleftharpoons}
 {\scriptstyle K_{#2}}}~}
\makeatletter
\def\th@plain{%
\thm@notefont{}% same as heading font
\itshape % body font
}
\def\th@definition{%
\thm@notefont{}% same as heading font
\normalfont % body font
}
\makeatother
\date{}


\title{Lecture-15: $L^p$ convergence}% of random variables}
\author{}

\begin{document}
\maketitle


%\section{$L^p$ space}
%
%\begin{defn}[$L^p$ space] 
%Consider a probability space $(\Omega, \sF, P)$. 
%For any $p > 1$, we say that a random variable $X \in L^p$, if $(\E\abs{X}^p)^{\frac{1}{p}} < \infty$, and we can define a norm 
%\EQ{
%\norm{X}_p \triangleq (\E\abs{X}^p)^{\frac{1}{p}}.
%}
%\end{defn}
%%\begin{thm}[Minkowski's inequality] 
%%Norm on the $L^p$ satisfies the triangle inequality. 
%%That is, if $X, Y \in L^p$, then $X+Y\in L^p$ and 
%%\EQ{
%%\norm{X+Y}_p \le \norm{X}_p + \norm{Y}_p.
%%}
%%\end{thm}
%%\begin{proof} 
%%We first show that $X+Y\in L^p$. 
%%To this end, it suffices to show that $\norm{X+Y}_p$ is bounded if $\norm{X}_p$ and $\norm{Y}_p$ are bounded.  
%%From the convexity of $g(x) = \abs{x}^p$ for $p \ge 1$, we get 
%%\EQ{
%%\abs{X+Y}^p \le \frac{1}{2}\abs{2X}^p +\frac{1}{2}\abs{2Y}^p = 2^{p-1}(\abs{X}^p +\abs{Y}^p).
%%}
%%Taking expectation on both sides, it follows from the linearity of expectation, $\norm{X+Y}_p \le  2^{p-1}(\norm{X}_p + \norm{Y}_p)$. 
%%
%%We now show that $\norm{}_p$ is a norm and satisfies the triangle inequality. 
%%From the triangle equality $\abs{X+Y} \le \abs{X} + \abs{Y}$ for two random variables $X,Y \in L^p$,  the linearity of expectation, and the H\"{o}lder inequality for pair of random variables $X,  (X+Y)^{p-1}$ and $Y, (X+Y)^{p-1}$, we get 
%%\EQ{
%%\E\abs{X+Y}^p \le \E\abs{X}\abs{X+Y}^{p-1} + \E\abs{Y}\abs{X+Y}^{p-1} \le ((\E\abs{X}^p)^{\frac{1}{p}}+ (\E\abs{Y}^p)^{\frac{1}{p}})(\E\abs{X+Y}^p)^{1-\frac{1}{p}}.
%%}
%%\end{proof}
%%%\begin{shaded*}
%%\begin{thm}
%%For a probability space $(\Omega,\sF, P)$ and $q \ge p \ge 1$, 
%%we have $L^q\subseteq L^p$.  
%%\end{thm}
%%\begin{proof} 
%%Consider $q \ge p \ge 1$, and a random variable $X \in L^q$ defined on the probability space $(\Omega, \sF, P)$. 
%%Applying H\"{o}lder's inequality to the product of random variables $\abs{X}^{p}\cdot 1$ with conjugate variables $p' \triangleq \frac{q}{p} \ge 1$ and $q' \triangleq \frac{q}{q-p} \ge 1$, 
%%we get %the result 
%%%\EQ{
%%$\E\abs{X}^p = \E[\abs{X}^{\frac{q}{p'}}\cdot 1] \le (\E\abs{X}^q)^{\frac{1}{p'}}.$
%%%}
%%\end{proof}
%%%\end{shaded*}
%\begin{shaded*}
%\begin{exmp}[Mean square error] 
%Consider a sequence of random variables $X: \Omega \to \R^\N$ such that 
%\begin{xalignat*}{2}
%&m \triangleq \E X_n,&
%&\rho_k \triangleq \cov(X_nX_{n+k})\text{ for all }n,k \in \N.
%\end{xalignat*}
%The \textbf{best affine predictor} of $X_{n+1}$ based on $X_1, \dots, X_n$ is given by $\hat{X}_{n+1} = \alpha_0 + \sum_{i=1}^n\alpha_iX_i$ for  $(\alpha_0, \alpha_1, \dots, \alpha_n) \in \R^{n+1}$ such that the \textbf{mean square error} is minimized. 
%
%Taking $\alpha_{n+1} = -1$, we can write the mean square error as 
%\EQ{
%\E\abs{\hat{X}_{n+1}-X_{n+1}}^2 
%= \E\abs{\alpha_0 +\sum_{i=1}^{n+1}\alpha_im+\sum_{i=1}^{n+1}\alpha_i(X_i-m)}^2 
%= (\alpha_0 +\sum_{i=1}^{n+1}\alpha_im)^2 + \sum_{i=1}^{n+1}\sum_{j=1}^{n+1}\alpha_i\alpha_j\rho_{\abs{j-i}}.
%}
%Taking derivative with respect to coefficient $\alpha_0$ and equating it to zero, we get $\alpha_0 = -m + m\sum_{i=1}^{n}\alpha_i$. 
%Taking derivatives with respect to  coefficients $\alpha \in \R^{n}$ and equating it to zero, 
%we get 
%\EQ{
%\sum_{j=1}^{n}\alpha_j\rho_{\abs{j-i}}  = \rho_{n+1-i},~\quad i \in [n].
%}

%Taking $\alpha_0 = -1$, 
%we have
%\EQ{
%\E\abs{X_{n+1}-\hat{X}_{n+1}}^2 
%%= \min_{\alpha \in \R^n}\E(X_{n+1}-\sum_{i=1}^n\alpha_iX_i)^2
%= \min_{\alpha \in \R^n}\Big(\rho_0 + m^2 + \sum_{i=1}^n\alpha_i^2(\rho_0+m^2) 
%- 2\sum_{i=1}^n\alpha_i((m^2+\rho_{n+1-i})-\sum_{k=1}^{n-i}\alpha_{i+k}(m^2+\rho_k))\Big).
%}
%Taking derivatives with respect to coefficients $\alpha \in \R^n$, we get 
%\EQ{
%\alpha_i(\rho_0+m^2) = m^2+\rho_{n+1-i}- \sum_{k=1}^{n-i}\alpha_{i+k}(m^2+\rho_k),\quad i \in [n].
%}
%%This is achieved for the coefficients $\alpha \in \R^n$, such that
%%\EQ{
%%\alpha_i\E[X_i^2] = \E[X_iX_{n+1}],~i \in [n].
%%}
%%That is , $\alpha_i = \frac{\rho_{n+1-i}+m^2}{\rho_0+m^2}$ for all $i \in [n]$. 
%\end{exmp}
%\end{shaded*}

\section{$L^p$ convergence}
\begin{defn}[Convergence in $L^p$] 
Let $p \ge 1$, 
then we say that a random sequence $X: \Omega \to \R^\N$ defined on a probability space $(\Omega,\sF,P)$ \textbf{converges in $L^p$} to a random variable $X_\infty:\Omega \to \R$, 
if 
\EQ{
\lim_n\norm{X_n - X_\infty}_p = 0.
}
The convergence in $L^p$ is denoted by $\lim_nX_n = X_\infty$ in $L^p$. 
\end{defn}
\begin{rem}
For $p \in [1, \infty)$, the convergence of a random sequence $X: \Omega\to\R^\N$ in $L^p$  to a random variable $X_\infty:\Omega \to \R$ is equivalent to 
\EQ{
\lim_n\E\abs{X_n-X_\infty}^p = 0.
}
\end{rem}
\begin{prop}[Convergences $L^p$ implies in probability] 
Consider $p \in [1, \infty)$ and a sequence of random variables $X:\Omega\to\R^\N$ defined on a probability space $(\Omega,\sF, P)$ such that $\lim_nX_n = X_\infty$ in $L^p$, then $\lim_nX_n = X_\infty$ in probability.
\end{prop}
\begin{proof}
Let $\epsilon > 0$, then from the Markov's inequality applied to random variable $\abs{X_n-X}^p$, we have
\EQ{
P\set{\abs{X_n-X_\infty} > \epsilon} \le \frac{\E\abs{X_n-X_\infty}^p}{\epsilon}.
}
\end{proof}
\begin{shaded*}
\begin{exmp}[Convergence almost surely doesn't imply convergence in $L^p$] 
Consider the probability space $([0,1], \sB([0,1]), \lambda)$ such that $\lambda([a,b]) = b-a$ for all $0 \le a \le b \le 1$. 
We define the scaled indicator random variable $X_n: \Omega \to \set{0,1}$ such that
\EQ{
X_n(\omega) = 2^n\Ind{\left[0, \frac{1}{n}\right]}(\omega).
}
We define $N = \set{0}$, 
and for any $\omega \notin N$, we can find $m \triangleq \lceil \frac{1}{\omega}\rceil$, 
such that for all $n > m$, we have $X_n(\omega) = 0$. 
Since $\lambda(N) = 0$, it implies that $\lim_nX_n = 0$ a.s. 
%Then, $\lim_nX_n = 0$ in probability, since for any $1 > \epsilon > 0$, we have 
%\EQ{
%P\set{\abs{X_n} > \epsilon} = \frac{1}{n}. 
%}
However, we see that $\E\abs{X_n}^p = \frac{2^{np}}{n}$. 
\end{exmp}
\end{shaded*}
\begin{rem}
Convergence almost surely implies convergence in probability. 
Therefore, above example also serves as a counterexample to the fact that convergence in probability doesn't imply convergence in $L^p$.
\end{rem}
%\begin{shaded*}
\begin{thm}[$L^2$ weak law of large numbers] 
Consider a sequence of uncorrelated random variables $X:\Omega\to\R^\N$ defined on a probability space $(\Omega,\sF, P)$ such that $\E X_n = \mu$ and $\Var(X_n) = \sigma^2$ for all $n \in \N$. 
Defining the sum $S_n\triangleq \sum_{i=1}^nX_i$ and the $n$-empirical mean $\bar{X}_n \triangleq \frac{S_n}{n}$, 
we have $\lim_n\bar{X}_n=\mu$ in $L^2$ and in probability.
\end{thm}
\begin{proof} 
From the uncorrelatedness of random sequence $X$, 
and linearity of expectation, we get
\EQ{
\Var(\bar{X}_n) = \E(\bar{X}_n-\mu)^2 = \frac{1}{n^2}\E(S_n-n\mu)^2 = \frac{\sigma^2}{n}.
}
It follows that $\lim_n\bar{X}_n = \mu$ in $L^2$. 
Since the convergence in $L^p$ implies convergence in probability, the result holds.
\end{proof}
\begin{thm}[$L^1$ weak law of large numbers] 
Consider an \iid random sequence $X:\Omega\to\R^\N$ defined on a probability space $(\Omega,\sF, P)$ such that $\E\abs{X_1}< \infty$ and $\E X_1 = \mu$. 
Defining the sum $S_n\triangleq \sum_{i=1}^nX_i$ and the $n$-empirical mean $\bar{X}_n \triangleq \frac{S_n}{n}$, 
we have $\lim_n\bar{X}_n=\mu$ in probability.
\end{thm}
%\end{shaded*}
\begin{shaded*}
\begin{exmp}[Convergence in $L^p$ doesn't imply almost surely] 
Consider the probability space $([0,1], \sB([0,1]), \lambda)$ such that $\lambda([a,b]) = b-a$ for all $0 \le a \le b \le 1$. 
For each $k \in \N$, we consider the sequence $S_k = \sum_{i=1}^ki$, and define integer intervals $I_k \triangleq \set{S_{k-1}+1, \dots, S_{k}}$. 
Clearly, the intervals $(I_k:k \in \N)$ partition the natural numbers, 
and each $n \in \N$ lies in some $I_{k_n}$, such that $n = S_{k+n-1}+i_n$ for $i_n \in [k_n]$. 
Therefore, for each $n \in \N$, we define indicator random variable $X_n: \Omega \to \set{0,1}$ such that 
\EQ{
X_n(\omega) \triangleq \Ind{\left[\frac{i_n-1}{k_n}, \frac{i_n}{k_n}\right]}(\omega).
}
For any $\omega \in [0,1]$, we have $X_n(\omega) = 1$ for infinitely many values since there exist infinitely many $(i,k)$ pairs such that $\frac{(i-1)}{k} \le \omega \le \frac{i}{k}$, 
and hence $\lim\sup_nX_n(\omega) = 1$ and hence $\lim_nX_n(\omega) \neq 0$. 
However, $\lim_nX_n(\omega) = 0$ in $L^p$, since 
\EQ{
\E\abs{X_n}^p = \lambda\set{X_n(\omega) \neq 0} =\frac{1}{k_n} .
}
\end{exmp}
\end{shaded*}
%\begin{thm}[$L^2$ weak law of large numbers] 
%Let $(X_n: n \in \N)$ be a random sequence of uncorrelated random variables with $\E X_n = \mu$ and $\Var(X_i) 
%\le C < \infty$. 
%If $S_n = \sum_{i=1}^nX_i$ then $$ in L2 and in probability.
%\end{thm}


\section{$L^1$ convergence theorems}
\begin{thm}[Monotone Convergence Theorem]
Consider a non-decreasing non-negative random sequence $X: \Omega\to\R_+^\N$ defined on a probability space $(\Omega,\sF,P)$, 
such that $X_n \in L^1$ for all $n \in \N$. 
Let $X_\infty(\omega) \triangleq \sup_nX_n(\omega)$ for all $\omega \in \Omega$, then  
%\EQ{
$
\E X_\infty = \sup_n\E X_n.
$
%}
\end{thm}
\begin{proof}
From the monotonicity of sequence $X$ and the monotonicity of expectation, we have $\sup_n\E X_n \le \E X_\infty$. 
Let $\alpha \in (0,1)$ and $Y: \Omega\to\R_+$ a non-negative simple random variable such that $Y \le X_\infty$. 
We define 
\EQ{
E_n \triangleq \set{\omega \in \Omega: X_n(\omega) \ge \alpha Y} \in \sF. 
} 
From the monotonicity of sequence $X$, 
the sequence of events $E \in \sF^\N$ are monotonically non-decreasing such that $\cup_{n \in \N}E_n = \Omega$. 
It follows that 
\EQ{
\alpha\E[Y\Ind{E_n}] \le \E[X_n\Ind{E_n}] \le \E X_n.
}
We will use the fact that $\lim_n\E[Y\Ind{E_n}] = \E[Y]$, then $\alpha\E Y \le \sup_n\E X_n$. 
Taking supremum over all $\alpha \in (0,1)$ and all simple functions $Y \le X_\infty$, we get $\E X_\infty \le \sup_n\E X_n$. 
\end{proof}
\begin{thm}[Fatou's Lemma]
Consider a non-negative random sequence $X: \Omega\to\R_+^\N$ defined on a probability space $(\Omega,\sF,P)$. 
Let $X_\infty(\omega) \triangleq \lim\inf_nX_n(\omega)$ for all $\omega\in\Omega$, then 
%\EQ{
$
\E X_\infty \le \lim\inf_n\E X_n.
$
%} 
\end{thm}
\begin{proof}
We define $Y_n \triangleq \inf_{k \ge n}X_k$ for all $n \in \N$. 
It follows that $Y: \Omega \to \R_+^\N$ is a non-negative non-decreasing sequence of random variables, 
and $X_\infty = \sup_nY_n = \lim_nY_n$. 
Applying monotone convergence theorem to random sequence $Y$, we get $\E X_\infty = \sup_n\E Y_n$. 
The result follows from the monotonicity of expectation, 
and the fact that $Y_n \le X_k$ for all $k \ge n$, to get $\E Y_n \le \inf_{k\ge n}\E X_k$.  
\end{proof}
\begin{thm}[Dominated Convergence Theorem]
Let $X: \Omega \to \R^\N$ be a random sequence defined on a probability space $(\Omega, \sF, P)$. %, such that $X_n \in L^1$ for all $n \in \N$. 
If $\lim_nX_n = X_\infty$ a.s. and there exists a $Y: \Omega \to \R_+$ such that $Y \in L^1$ and $\abs{X_n} \le Y$ a.s. for all $n\in\N$, 
then $\E X_\infty = \lim_n\E X_n$. 
\end{thm}
\begin{proof} 
From the hypothesis, we have $Y+X_n \ge 0$ a.s. and $Y-X_n \ge 0$ a.s. 
Therefore, from Fatou's Lemma and linearity of expectation, we have 
\begin{xalignat*}{2}
&\E Y + \E X_\infty \le \liminf_n \E(Y+X_n) = \E Y + \liminf_n\E X_n,&
&\E Y - \E X_\infty \le \liminf_n \E(Y-X_n) = \E Y - \limsup_n\E X_n. 
\end{xalignat*} 
Therefore, we have $\limsup_n\E X_n \le \E X_\infty \le \liminf_n\E X_n$, and the result follows. 
\end{proof}

\section{Convergence theorems for conditional means}

\begin{prop} 
Let $X:\Omega\to\R^\N$ be a random sequence on the probability space $(\Omega, \sF, P)$ such that $\E\abs{X_n} < \infty$ for all $n \in \N$.  
Let $\sG$ and $\sH$ be event spaces such that $\sG, \sH \subset \sF$. 
Then the following theorems hold.
\begin{enumerate}
\item \textbf{Conditional monotone convergence theorem:} 
%If $0 \le X_n \le X_{n+1}$ a.s., for 
Consider a random sequence $X:\Omega\to\R_+^\N$ that is non-negative and non-decreasing a.s..
We define $X_\infty:\Omega\to\R_+$ for each $\omega \in \Omega$ as $X_\infty(\omega)\triangleq \sup_{n \in \N}X_n(\omega)$.  
If $X_\infty \in L^1$, then $\E[X_n\mid \sG] \uparrow \E[X_\infty\mid\sG]$ a.s. 

\item \textbf{Conditional Fatou's lemma:} 
Consider a random sequence $X:\Omega\to\R_+^\N$ that is non-negative and non-decreasing a.s.. 
If $\lim\inf_n X_n \in L^1$, then $\E[\lim\inf_nX_n\mid \sG] \le \lim\inf_n\E[X_n\mid \sG]$ a.s. 

\item \textbf{Conditional dominated convergence theorem:} 
Consider a random sequence $X:\Omega\to\R_+^\N$ and a random variable $Y:\Omega\to\R_+$ such that $Y \in L^1$ and $\abs{X_n} \le Y$ a.s. for all $n \in  \N$. 
If $\lim_n X_n = X_\infty$ a.s., then $\E[X_n\mid \sG] \to  \E[X_\infty\mid\sG]$ a.s. and in $L^1$.
\end{enumerate}
\end{prop}
\begin{proof}
Let $X:\Omega\to\R^\N$ be a random sequence on the probability space $(\Omega, \sF, P)$ such that $X_n \in L^1$ for all $n \in \N$.  
%Note: Some of the properties are proved in detail. 
%The others are only commented upon, since they are either similar to the other ones or otherwise not hard. 
\begin{enumerate}
\item \textbf{Conditional monotone-convergence theorem:} 
By monotonicity, we have $\E[X_n\mid \sG] \uparrow Y$ a.s. where $Y:\Omega\to \R_+$ is $\sG$ measurable. 
The monotone convergence theorem implies that, for each $G \in  \sG$,
\EQ{
\E[Y \Ind{G}] = \lim_n
\E[\Ind{G}\E[X_n\mid \sG]] = \lim_n
\E[\Ind{G}X_n] = \E[\Ind{G}X_\infty].
}

\item \textbf{Conditional Fatou's lemma:} 
Defining $Y_n \triangleq \inf_{k\ge n} X_k$, we get $Y_n \uparrow Y_\infty = \lim\inf_k X_k$. 
By monotonicity,
\EQ{
\E[Y_n\mid \sG] \le \inf_{k\ge n}\E[X_k\mid \sG]\text{ a.s.}.
}
The conditional monotone-convergence theorem implies that
\EQ{
\E[Y_\infty\mid \sG] = \lim_{n\in \N}\E[Y_n\mid \sG] \le \lim\inf_n\E[X_n\mid \sG]\text{ a.s.}.
}
\item \textbf{Conditional dominated-convergence theorem:}
By the conditional Fatou's lemma, we have
\meq{2}{
&\E[Z + X_\infty\mid \sG] \le \lim\inf_n\E[Z + X_n\mid \sG]\text{ a.s. },&&
%\text{ as well as  }
\E[Z - X_\infty\mid \sG] \le \lim\inf_n\E[Z - X_n\mid \sG] \text{ a.s. },
}
and the a.s.-statement follows.
\end{enumerate}
\end{proof}

\section{Uniform integrability}
\red{
\begin{defn}[uniform integrability] 
A family $X_t \in (L^1)^T$ of random variables indexed by $T$ is \textbf{uniformly integrable} if 
\EQ{
\lim_{a \to \infty}\sup_{t \in T}\E[\abs{X_t}\SetIn{\abs{X_t}> a}] = 0.
}
\end{defn}
\begin{shaded*}
\begin{exmp}[Single element family] 
If $\abs{T} = 1$, then the family is uniformly integrable, 
since $X_1 \in L^1$ and $\lim_a\E[\abs{X_1}\SetIn{\abs{X_t} > a}] = 0$. 
This is due to the fact that $(X_n \triangleq \abs{X}\SetIn{\abs{X} \le n}: n \in \N)$ is a sequence of increasing random variables $\lim_nX_n  = X$. 
From monotone convergence theorem, we get  $\lim_n\E\abs{X_n} = \E\lim_n\abs{X_n}$. 
Therefore, 
\EQ{
\lim_a\E[\abs{X}\SetIn{\abs{X} > a}] = \E\abs{X} - \lim_a\E[\abs{X}\SetIn{\abs{X} \le a}] = 0. 
}
\end{exmp}
\end{shaded*}
\begin{prop}
Let $X \in L^p$ and $(A_n: n \in \N) \subset \sF$ be a sequence of events such that $\lim_nP(A_n) = 0$, then 
\EQ{
\lim_n\norm{\abs{X}\mathbbm{1}_{A_n}}_p = 0.
}
\end{prop}
\begin{shaded*}
\begin{exmp}[Dominated family] 
If there exists $Y \in L^1$ such that $\sup_{t \in T}\abs{X_t} \le \abs{Y}$, 
then the family of random variables $(X_t: t \in T)$ is uniformly integrable. 
This is due to the fact that 
\EQ{
\sup_{t \in T}\E[\abs{X}\SetIn{\abs{X} > a}] \le \E[\abs{Y}\SetIn{\abs{Y} > a}]. 
}
\end{exmp}
\end{shaded*}
\begin{shaded*}
\begin{exmp}[Finite family] then the family of random variables $(X_t: t \in T)$ is uniformly integrable. 
This is due to the fact that $\sup_{t \in T}\abs{X_t} \le \sum_{t \in T}\abs{X_t} \in L^1$. 
\end{exmp}
\end{shaded*}
\begin{thm}[Convergence in probability with uniform integrability implies convergence in $L^p$] 
Consider a sequence of random variables $(X_n: n\in \N) \subset L^p$ for $p \ge 1$.  
Then the following are equivalent. 
\begin{enumerate}[(a)]
\item The  sequence $(X_n: n\in \N)$ converges in $L^p$, i.e. $\lim_n\E\abs{X_n -  X}^p = 0$.
\item The sequence  $(X_n: n\in \N)$ is Cauchy in $L^p$, i.e. $\lim_{m,n \to \infty}\E\abs{X_n-X_m}^p = 0$. 
\item $\lim_nX_n = X$ in probability and the sequence $(\abs{X_n}^p: n \in \N)$ is uniformly integrable.  
\end{enumerate}
\end{thm}
\begin{proof} 
For a random sequence $(X_n: n\in \N)$ in $L^p$, 
we will show that $(a)\implies (b) \implies (c) \implies (a)$. 
\begin{enumerate}[(a)]
\item[$(a)\implies (b):$]  We assume the sequence $(X_n: n\in \N)$ converges in $L^p$. 
Then, from Minkowski's inequality, we can write 
\EQ{
 (\E\abs{X_n-X_m}^p)^{\frac{1}{p}} \le (\E\abs{X_n-X}^p)^{\frac{1}{p}} + (\E\abs{X_m-X}^p)^{\frac{1}{p}}.
}
\item[$(b) \implies (c):$] We assume that the sequence  $(X_n: n\in \N)$ is Cauchy in $L^p$, i.e. $\lim_{m,n \to \infty}\E\abs{X_n-X_m}^p = 0$.  
Let $\epsilon > 0$, then for each $n \in \N$, there exists $N_\epsilon$ such that for all $n,m \ge N_\epsilon$
\EQ{
\E\abs{X_n-X_m}^p \le \frac{\epsilon}{2}.
}
Let $A_a =\set{\omega \in A: \abs{X_n} > a}$. 
Then, using triangle inequality and the fact that $\mathbbm{1}_{A_a} \le 1$, from the linearity and monotonicity of expectation, 
we can write for $n \ge N_\epsilon$ 
\EQ{
(\E[\abs{X_n}^p\SetIn{\abs{X_n} > a}])^{\frac{1}{p}}
\le (\E[\abs{X_{N_\epsilon}}^p\mathbbm{1}_{A_a}])^{\frac{1}{p}}  + (\E[\abs{X_n - X_{N_\epsilon}}^p] )^{\frac{1}{p}}
\le (\E[\abs{X_{N_\epsilon}}^p\mathbbm{1}_{A_a}])^{\frac{1}{p}}+ \frac{\epsilon}{2}.
}
Therefore, we can write $\sup_n \E[\abs{X_n}^p\SetIn{\abs{X_n} > a}]  \le \sup_{m \le N_\epsilon}\E[\abs{X_m}^p\mathbbm{1}_{A_a}]+ \frac{\epsilon}{2}$. 
Since $(\abs{X_n}^p: n\le N_\epsilon)$ is finite family of random variables in $L^1$, it is uniformly integrable. 
Therefore, there exists $a_\epsilon \in \R_+$ such that $\sup_{m \le N_\epsilon}(\E[\abs{X_m}^p\mathbbm{1}_{A_a}])^{\frac{1}{p}} < \frac{\epsilon}{2}$. 
Taking $a' = \max\set{a, a_\epsilon}$, we get $\sup_n (\E[\abs{X_n}^p\SetIn{\abs{X_n} > a'}])^{\frac{1}{p}} \le \epsilon$. 
Since the choice of $\epsilon$ was arbitrary, it follows that 
\EQ{
\lim_{a \to \infty}\sup_n (\E[\abs{X_n}^p\SetIn{\abs{X_n} > a'}])^{\frac{1}{p}} = 0.
}
The convergence in probability follows from the Markov inequality, i.e.
\EQ{
P\set{\abs{X_n-X_m}^p > \epsilon} \le \frac{1}{\epsilon}\E\abs{X_n-X_m}^p. 
}
\item[$(c)\implies (a):$] 
%Let $\lim_nX_n = X$ in probability and the sequence $(\abs{X_n}^p: n \in \N)$ be uniformly integrable.  
Since the sequence $(X_n: n\in \N)$ is convergent in probability to a random variable $X$, there exists a subsequence $(n_k: k \in \N) \subset \N$ such that $\lim_kX_{n_k} = X$ a.s. 
Since $(\abs{X_n}^p: n \in \N)$ is a family of uniformly integrable sequence, by Fatou's Lemma
\EQ{
\E\abs{X}^p \le \lim\inf_k\E\abs{X_{n_k}}^p \le \sup_n\E\abs{X_n}^p < \infty.
}
Therefore, $X \in L^1$, and we define $A_n(\epsilon) = \set{\abs{X_n-X} > \epsilon}$ for any $\epsilon > 0$. 
From Minkowski's inequality, we get 
\EQ{
\norm{X_n-X}_p \le \norm{(X_n-X)\SetIn{\abs{X_n-X}^p \le \epsilon}}_p+\norm{X_n\mathbbm{1}_{A_n(\epsilon)}}_p+\norm{X\mathbbm{1}_{A_n(\epsilon)}}_p.
}
We can check that $\norm{(X_n-X)\mathbbm{1}_{A_n^c(\epsilon)}}_p \le \epsilon$. 
Further, since $\lim_nX_n = X$ in probability, $(A_n: n \in \N) \subset \sF$ is decreasing sequence of events, 
and since $X_n, X \in L^1$, we have $\lim_n\norm{X_n\mathbbm{1}_{A_n(\epsilon)}} = \lim_n\norm{X\mathbbm{1}_{A_n(\epsilon)}} = 0$.
%Then, for any $\epsilon > 0$, we can find $n_\epsilon$ such that for all $n \ge n_\epsilon$
%\EQ{
%P\set{\omega \in \Omega: \abs{X_n-X}(\omega) > \epsilon} 
%}
\end{enumerate}
\end{proof}
}
\end{document}